\SourceLeaf{226}%
\noindent Now two branches of each great circle issue from the point \(A\), so the two angles formed there add up to \(180^0\). Our analysis shows that we must take the branches whose directions agree with the sense in which we proceed from \(L\) to \(L'\) and from \(L''\) to \(L'''\). It is then clear that, since the great circles intersect at two points, we may choose either one. In place of the angle \(A\) we may also use the arc between the poles of the great circles of which \(LL',L''L'''\) are arcs. Evidently we must choose poles that lie on the same side relative to their respective arcs. For instance, as we proceed from \(L\) toward \(L'\) and from \(L''\) toward \(L'''\), both poles may lie on the right or both on the left.

VII. Let \(L,L',L''\) be three points on the spherical surface, and for brevity set
\[
\begin{aligned}
\cos(1)L&=x,&\cos(2)L&=y,&\cos(3)L&=z\\
\cos(1)L'&=x',&\cos(2)L'&=y',&\cos(3)L'&=z'\\
\cos(1)L''&=x'',&\cos(2)L''&=y'',&\cos(3)L''&=z''
\end{aligned}
\]
\noindent and also
\[
xy'z''+x'y''z+x''yz'-xy''z'-x'yz''-x''y'z=\Delta
\]
\noindent Let \(\lambda\) denote the pole of the great circle containing the arc \(LL'\), choosing the pole that lies in the same position relative to this arc as the point \((1)\) does relative to the arc \((2)(3)\). By the preceding theorem, \(yz'-y'z=\cos(1)\lambda\mathord{.}\sin(2)(3)\mathord{.}\sin LL'\), or, since \((2)(3)=90^0\),
\[
\begin{aligned}
yz'-y'z&=\cos(1)\lambda\mathord{.}\sin LL',\ \text{and similarly}\\
zx'-z'x&=\cos(2)\lambda\mathord{.}\sin LL'\\
xy'-x'y&=\cos(3)\lambda\mathord{.}\sin LL'
\end{aligned}
\]
\noindent Multiplying these equations by \(x'',y'',z''\), respectively, and adding them, we obtain, with the help of the second theorem stated in \(V\),
\[
\Delta=\cos\lambda L''\mathord{.}\sin LL'
\]
\noindent We must now distinguish three cases. \textit{First}, whenever \(L''\) lies on the same great circle that contains the arc \(LL'\), we have \(\lambda L''=90^0\) and therefore \(\Delta=0\). When \(L''\) lies off that great circle, we have the \textit{second} case if it is on the same side as \(\lambda\), and the \textit{third} if it is on the opposite side. In these two cases the points \(L,L',L''\)
